Joined BGP and active-probing/network-telescope measurement, applied to
Iran's 2025--26 shutdown at 8-hour resolution across the full period and at
BGP-update resolution across five historical event windows, supports a
filtering-dominant mechanism (H1): routes stayed announced while active
reachability collapsed, a pattern absent from a matched control population
and corroborated by an independently authored, methodologically distinct
study of the same event. Restoration selectivity (H3) is real but coarser
than a clean per-sector ranking -- government and core state infrastructure
recover fastest and most completely, the unclassified long tail slowest and
least completely, with genuine complications (mobile's fast typical delay
but long tail of stragglers) stated outright, not smoothed away. H2's specific
topological claim -- gateway-level versus distributed disconnection -- and
H4's lack of a cross-window control check are identified as open
limitations, not silently folded into the headline results.

Every number in this paper can be traced back to a specific, documented
analytical choice, and every threshold and robustness check was fixed
before its effect on a result was known. That discipline caught two real
mistakes during this study, and both are reported here rather than quietly
corrected and forgotten: an early version of the address-space calculation
double-counted overlapping BGP announcements, and the threshold separating
``reachable'' from ``dark'' was left at a placeholder value for longer than it
should have been before being properly calibrated. Neither changed this
paper's main conclusions, but both are the kind of error that a less
disciplined analysis could easily have missed.

Extending H2's topological decomposition to test the
gateway-versus-access-network claim directly, and building an H4-side
control-population check comparable to H1/H3's, are the most direct next
steps this study's own limitations point to.
